\documentclass[a4paper,10pt]{article}

\usepackage{fontspec}
\usepackage{xeCJK}
\usepackage{parskip}
\usepackage[table,svgnames]{xcolor}
\usepackage{geometry}
\usepackage{hyperref}
\usepackage{titlesec}
\usepackage{tabularx}
\usepackage{array}
\usepackage{booktabs}
\usepackage{enumitem}
\usepackage{multicol}
\usepackage{fontawesome5}
\usepackage{tikz}
\usepackage{ifthen}

\defaultfontfeatures{Mapping=tex-text}

\setmainfont{texgyrepagella}[
    Extension=.otf,
    UprightFont=*-regular,
    BoldFont=*-bold,
    ItalicFont=*-italic,
    BoldItalicFont=*-bolditalic
]
\setCJKmainfont{Noto Serif CJK SC}
\setCJKsansfont{Noto Sans CJK SC}

\geometry{left=2cm,right=2cm,top=2.2cm,bottom=2.2cm}
\pagestyle{empty}

\definecolor{themecolor}{RGB}{30,65,125}
\definecolor{lightline}{RGB}{190,205,225}
\definecolor{accentbg}{RGB}{240,245,252}
\definecolor{tagbg}{RGB}{232,240,250}
\definecolor{tagtext}{RGB}{40,70,130}
\definecolor{lightgray}{RGB}{120,120,120}
\definecolor{skillbar}{RGB}{30,65,125}
\definecolor{skillbarbg}{RGB}{225,232,242}

\hypersetup{
    colorlinks=true,
    breaklinks=true,
    urlcolor=themecolor,
    linkcolor=themecolor
}


\titleformat{\section}{\large\bfseries\color{themecolor}}{}{0em}{}[{\color{lightline}\titlerule}]
\titlespacing{\section}{0pt}{2em}{1em}

\newcommand{\skillbar}[2]{%
    \begin{tikzpicture}[baseline=-0.3em]
        \fill[skillbarbg] (0,0) rectangle (3.5em,0.35em);
        \fill[skillbar] (0,0) rectangle ({3.5*#2/5}em,0.35em);
    \end{tikzpicture}%
}

\newcommand{\cvtag}[1]{%
    \tikz[baseline=-0.5ex]\node[fill=tagbg, text=tagtext, rounded corners=2pt, inner xsep=4pt, inner ysep=1.5pt, font=\small\bfseries]{#1};%
}

\newcommand{\contactitem}[2]{\raisebox{-0.15ex}{#1}\,#2}

\begin{document}

\begin{center}
    {\Huge\bfseries\color{themecolor} 黄煜珩}\\[6pt]
    {\large\color{lightgray} 宁波大学\,\textbar\,计算机科学与技术\,\textbar\,学硕}\\[8pt]
    {\small\color{themecolor}
    \contactitem{\faPhone}{139\,6232\,0052} \quad
    \contactitem{\faEnvelope}{\href{mailto:2130588669@qq.com}{2130588669@qq.com}} \quad
    \contactitem{\faEnvelope}{\href{mailto:hihyhenry@gmail.com}{hihyhenry@gmail.com}}
    % \contactitem{\faGraduationCap}{\href{https://scholar.google.com/citations?user=JR3GvqAAAAAJ}{Google Scholar}} \quad
    % \contactitem{\faGithub}{\href{https://github.com/LahenVieLesry}{GitHub}} \quad
    % \contactitem{\faUniversity}{\href{https://www.nbu.edu.cn}{宁波大学}}
    }
\end{center}

\section{个人简介}

\begin{minipage}[t]{\textwidth}
\vspace{2pt}
宁波大学计算机科学与技术硕士应届生，专注\textbf{多媒体信息安全与人工智能安全}方向。硕士期间以第一作者在 \textbf{IEEE TIFS}（CCF-A / SCI 一区）投稿 2 篇、语音顶级会议\textbf{Interspeech}（Core A, CCF-B）发表 1 篇，累计参与发表/投稿论文 \textbf{10 篇}。具备从算法设计、模型训练到工程落地的全栈研发能力。
\vspace{8pt}
\end{minipage}

\begin{tabularx}{\textwidth}{@{}p{\textwidth}@{}}

\textbf{\color{themecolor}\faShieldAlt\;模型安全性评估研究} \\[3pt]
\parbox[t]{\textwidth}{\small
\textbullet\;\textbf{攻击侧}：设计梯度与特征空间样本筛选框架 + 中心覆盖多折选择算法，投毒率 $\leq 0.05\%$ 下攻击成功率从近 0 提升至 \textbf{70\%}\\
\textbullet\;开发基于音频压缩的后门攻击工具链 CBA，5\% 投毒率下达攻击成功率 $>$ \textbf{99\%}，发表于 \cvtag{Interspeech 2025}\\
\textbullet\;\textbf{防御侧}：构建动量孪生范式 MoSia 无数据后门防御框架，已在 \textbf{4 类任务、6 种攻击方法}上完成端到端验证
} & \\

\textbf{\color{themecolor}\faWaveSquare\;深度学习与音频信号处理} \\[3pt]
\parbox[t]{\textwidth}{\small
\textbullet\;PyTorch 全栈工程：数据预处理 $\rightarrow$ 自定义 Dataset/DataLoader $\rightarrow$ 分布式训练 $\rightarrow$ TensorBoard 可视化\\
\textbullet\;音频编解码\cvtag{MP3}\cvtag{AAC}\cvtag{FLAC} \;|\; 时频分析\cvtag{STFT}\cvtag{Mel}\cvtag{DWT} \;|\; 语音处理 pipeline
} & \\

\textbf{\color{themecolor}\faRobot\;LLM \& Agent 工程} \\[3pt]
\parbox[t]{\textwidth}{\small
\textbullet\;独立交付 \href{https://github.com/LahenVieLesry/valveye}{\textbf{Valveye}}：\cvtag{LangChain}\cvtag{LangGraph} Supervisor + 4 Specialist 多 Agent 系统\\
\textbullet\;三源容错降级、23 区域异步并发比价 + 实时汇率转换；Rich CLI 交互系统与多渠道消息推送\\
\textbullet\;设计 12 个领域工具，实现意图分解、任务路由与 Agent Handoff 协作\\
\textbullet\;集成 ReAct 多步推理与 \cvtag{OpenViking} L0/L1/L2 三层记忆架构（token 利用率 \textbf{+90\%}）\\
\textbullet\;暴露 \cvtag{MCP} Server 接口，支持外部工具客户端（Claude Code 等）无缝调用
} & \\

\end{tabularx}

\section{教育背景与获奖}

\begin{tabularx}{\textwidth}{@{}p{2.8cm}X@{}}
\textbf{2023.09 -- 至今} & \textbf{宁波大学}（浙江宁波）\quad 计算机科学与技术\;\textbf{学术型硕士} \\[1pt]
& GPA \textbf{3.55}/4.0 \;\textbar\; 排名 \textbf{前 10\%}\\[2pt]
& \cvtag{二等奖学金$\times$2} \;\cvtag{华为杯数模国三} \\[6pt]

\textbf{2019.09 -- 2023.06} & \textbf{苏州工学院}（江苏苏州）\quad 计算机科学与技术\;\textbf{学士} \\[1pt]
& GPA \textbf{3.49}/4.0 \;\textbar\; 排名 \textbf{前 10\%} \\[2pt]
& \cvtag{二等奖学金$\times$3} \;\cvtag{蓝桥杯国二} \;\cvtag{蓝桥杯省一} \;\cvtag{数竞国三} \\
\end{tabularx}

\section{专业技能}

% \begin{multicols}{1}
\begin{itemize}
    \item \textbf{操作系统：}Linux（Ubuntu / Debian / Arch）、Windows
    \item \textbf{编程语言：}Python、C/C++、Shell、SQL
    \item \textbf{深度学习：}PyTorch（TorchVision / TorchAudio / DataParallel）、Librosa、TensorBoard、Hugging Face
    \item \textbf{Agent：}Claude Code、Openclaw、Hermes、LangChain、工具编排、记忆管理、ReAct
    \item \textbf{音频信号处理：}STFT / ISTFT、Mel 频谱、小波变换、心理声学模型、音频编解码
    \item \textbf{数据科学：}NumPy、Pandas、Scikit-learn、SciPy、Matplotlib
    \item \textbf{工程与平台：}\LaTeX{}、Git、GitHub、APScheduler、Rich CLI、Docker、CMake
    \item \textbf{开发工具：}VS Code、Jupyter Notebook、SSH、PyCharm、MySQL
    \item \textbf{语言能力：}CET-6（544 分）\;\textbar\; 日语 (本科双学位)
\end{itemize}
% \end{multicols}

\newpage
\section{研究开发经历}

\begin{tabularx}{\textwidth}{@{}p{2.8cm}X@{}}
\textbf{2023.09 -- 至今} & \textbf{深度学习后门攻防研究} \quad {\small\color{lightgray} 宁波大学多媒体信息安全实验室} \\
& \par\smallskip
\par\smallskip\hangindent=1em\textbullet\;独立设计并实现\textbf{后门样本筛选系统}：基于 PyTorch 提取训练过程中的梯度与特征空间信息，构建重要性评分函数，配合中心覆盖多折选择算法保证筛选多样性，代码涵盖数据预处理、模型训练、自动化实验与结果可视化全流程
\par\smallskip\hangindent=1em\textbullet\;开发\textbf{基于音频压缩的后门攻击工具链}：实现比特压缩与小波压缩模块、可学习掩码优化器、特征域插值融合 pipeline，在 Google Speech Commands 公开数据集上完成端到端验证，攻击成功率 $>$ \textbf{99\%}
\par\smallskip\hangindent=1em\textbullet\;构建\textbf{无数据后门防御框架}：基于 PyTorch 实现孪生模型训练流程，设计 EMA 动量更新与梯度自相关计算模块，支持多种攻击方法的自动化防御评估
\\
\\

\textbf{2025.03 -- 至今} & \textbf{Valveye — 智能对话式 Steam 游戏价格助手}\\
& {\small\color{lightgray} 技术栈：LangChain + LangGraph + OpenViking + OpenAI API + MCP + Python + aiohttp + aiosqlite} \\
& \par
\par\smallskip\hangindent=1em\textbullet\;基于 LangChain + LangGraph 构建\textbf{ Supervisor + 4 Specialist 多 Agent 架构}，设计 12 个领域工具并按 Agent 分组，实现意图分解、任务路由与 Agent 间 Handoff 协作
\par\smallskip\hangindent=1em\textbullet\;实现\textbf{ ReAct 多步推理}策略（需求理解 → 候选获取 → 深度调查 → 综合推理 → 个性化呈现），支持多轮对话记忆（AsyncSqliteSaver 异步持久化）与结构化流式事件输出
\par\smallskip\hangindent=1em\textbullet\;集成 OpenViking \textbf{L0/L1/L2 三层渐进式记忆架构}，实现跨会话自动召回与长期记忆捕获
\par\smallskip\hangindent=1em\textbullet\;设计\textbf{多数据源容错}机制（ITAD/SteamDB/CheapShark 三源降级），23 个区域异步并发比价（asyncio.gather）+ 实时汇率转换
\par\smallskip\hangindent=1em\textbullet\;构建\textbf{完整 CLI 交互系统}：Rich Markdown 渲染、Agent 思考过程实时流式展示、对话持久化与多格式导出（Markdown/JSON/HTML）
\par\smallskip\hangindent=1em\textbullet\;实现\textbf{定时价格监控}调度（APScheduler）与 7 渠道富文本消息推送（Email、Telegram、Discord、企业微信/飞书/钉钉、QQ）
\par\smallskip\hangindent=1em\textbullet\;暴露 \textbf{MCP Server} 标准接口，支持外部工具客户端无缝调用；实现 Steam 热门游戏实时查询（热销/新品/特惠/即将推出）
\\
\end{tabularx}

\section{科研成果}

\vspace{2pt}
\begin{enumerate}[leftmargin=*, after=\strut, label={[\arabic*]}]

    \item \textbf{CBA: Backdoor Attack on Deep Speech Classification via Audio Compression}\\
    {\small \textbf{一作} \quad CCF-B \quad 已发表 \quad Proc. Interspeech 2025}

    \item \textbf{Breaking the Symmetry: Momentum Siamese Paradigm for Data-Free Proactive Backdoor Defense}\\
    {\small \textbf{一作} \quad CCF-A / SCI 一区 \quad 在审 \quad IEEE TIFS}

    \item \textbf{Breaking Ultra-Low Poisoning Ratio Barrier: Sample Selection for Backdoor Boundary Deconstruction}\\
    {\small \textbf{一作} \quad CCF-A / SCI 一区 \quad 在审 \quad IEEE TIFS}

    \item \textbf{From Compression to Trigger: Inaudible Backdoor Attack for Keyword Recognition System}\\
    {\small \textbf{一作} \quad SCI 四区 \quad 已录用 \quad IJAACS, 2026}

    \item \textbf{Amplifying Discriminative Distortions: A Generative Latent Feature Reinforcement Framework for Audio Spoofing Detection}\\
    {\small 三作 \quad CCF-C / SCI 二区 \quad 已发表 \quad ESWA, 2025}

    \item \textbf{DyMEvalNet: Dynamic Text-Audio-Personalization Fusion for Multimodal Music Quality Assessment}\\
    {\small 三作 \quad CCF-C \quad 已发表 \quad Proc. ASRU 2025}

    \item \textbf{Fine-Grained Feature Learning for Audio Deepfake Detection via Auxiliary Source Tracing Task}\\
    {\small 四作 \quad CCF-A / SCI 一区 \quad 在审 \quad IEEE TIFS}

    \item \textbf{Ouroboros: Self-Referential Backdoor Attacks on Speech Enhancement via Clean Audio Triggers}\\
    {\small 二作 \quad CCF-B \quad 在审 \quad Proc. Interspeech 2026}

    \item \textbf{Benign Backdoor for Active Speaker Protection in One-Shot Voice Conversion}\\
    {\small 二作 \quad CCF-B \quad 在审 \quad Proc. Interspeech 2026}

    \item \textbf{CAQA-Net: Continual Audio Quality Assessment Across Speech and Music Domains}\\
    {\small 三作 \quad CCF-B \quad 在审 \quad Proc. Interspeech 2026}
\end{enumerate}
\end{document}
